\documentclass[11pt,twoside,a4paper]{article}

\usepackage{amsmath, amssymb, amsthm}
\usepackage{geometry}
\geometry{a4paper, margin=1in}
\usepackage{graphicx}
\usepackage{hyperref}
\usepackage{booktabs}
\usepackage{cite}

\title{Geometric Routing via Riemannian-Complex Transport: \\
Invariant Signals and the Universal Closure Framework}
\author{Lionel Charles (Casey) Allard IV \\ 
\small Independent Researcher \\
\small caseyallard@hotmail.com}
\date{\today}

\begin{document}

\maketitle

\begin{abstract}
Building on prior work demonstrating sublinear ($K^{0.572}$) compute reduction via fixed Hopf-base sector discretization, we propose that the observed efficiency of geometric routing is a low-order projection of a broader topological field computation. Standard transformer architectures rely on Euclidean linear operators that inadequately capture phase-dependent, directional, and hierarchical structures. We introduce a geometric routing framework that generalizes scalar computation into a complex-valued local transport law over a Riemannian manifold. By embedding states within a complexified tangent space ($T_xM \otimes \mathbb{C}$), information propagation follows geodesic-aware routing. We introduce Isotropic Closure ($\sigma$) as a universal metric for structural resolution, demonstrating that both optimal network representations and additive combinatorial structures (prime $k$-tuples) occupy the same hyperbolic state space governed by $e^2$ scaling laws. Finally, we empirically validate this architecture via a train-free geometric probe, identifying a Space Invariant Signal ($\Delta_{sep} = 1.333$) that confirms geometric separation is strictly tied to local harmonic phase structure, independent of ambient space dimensionality. 
\end{abstract}

\section{Introduction: Beyond Euclidean Routing}
Recent empirical evaluations of token embeddings on the unit hypersphere established that angular non-uniformity can be exploited to replace learned gating matrices with fixed geometric maps. Specifically, a fixed 4D Hopf fibration yields a sparse routing footprint scaling at $K^{0.572}$, outperforming 100D adaptive K-means while preserving norm-invariance \cite{allard2026angular}. 

While the Hopf-base discretization provides an effective engineering surrogate for gating, it raises a fundamental architectural question: why does a fixed topological map organize semantic space more efficiently than a dense, learned Euclidean projection? 

This paper argues that standard attention and routing mechanisms suffer from a reliance on Euclidean linear operators, which treat information spread as diffusion across a flat grid. We propose a formal transition to a Riemannian-Complex architecture, where computation operates as a fluid dynamic system, and information follows geodesic paths across a curved latent landscape.

\section{Riemannian-Complex Transport Dynamics}
To enforce structural constraints natively, we define computation on a smooth Riemannian manifold $(M, g)$. Each latent state $x_t$ evolves on $M$. 

\subsection{Complexified Tangent Space}
At each point $x_t$, we define a local representation in the tangent space $v_t \in T_{x_t}M$. To introduce directional and phase-aware structure—allowing signals to interfere constructively or destructively—we extend this into a complexified tangent bundle:
\begin{equation}
z_t \in T_{x_t}M \otimes \mathbb{C}
\end{equation}
with the decomposition $z_t = v_t^{(r)} + i v_t^{(i)}$.

\subsection{Local Transport and Manifold Update}
We define a local update operator that encodes phase evolution (rotation) and amplitude modulation (scaling). Scalar complex multiplication generalizes into a learned complex transport field:
\begin{equation}
\mathcal{U}_{x_t}(z) = W_r z + i W_i z
\end{equation}
where $W_r$ and $W_i$ are learned linear maps on $T_{x_t}M$. The updated state is then projected back onto the manifold via the exponential map:
\begin{equation}
x_{t+1} = \exp_{x_t}\big(\Re(\mathcal{U}_{x_t}(z_t))\big)
\end{equation}
This geodesic-aware routing enforces curvature boundaries, replacing unconstrained linear accumulation with topology-preserving transport.

\section{The Universal Closure Framework}
The efficacy of the previously observed geometric routing implies the existence of a universal optimization attractor. We define this attractor via Isotropic Closure ($\sigma$), which quantifies the collapse of anisotropy. For a normalized simplex state $p$, the Quadratic Closure is:
\begin{equation}
\sigma_Q(X) = 1 - \frac{\sum_{i=1}^n\left(p_i - \frac{1}{n}\right)^2}{1 - \frac{1}{n}}
\end{equation}
where $\sigma = 1$ denotes full isotropic closure and $\sigma = 0$ denotes maximal dimensional concentration.

\subsection{The Complex Bridge and Hyperbolic Depth}
We map this closure onto a complex half-orbit anchored at $2i$ (minimal closure) and terminating at $0$ (perfect resolution). Transforming this circular phase to the Poincar\'{e} upper half-plane yields the Tangent Operator, acting as the projective slope:
\begin{equation}
\mathcal{T}(X) = \tan\left(\frac{\pi}{2}\sigma(X)\right)
\end{equation}
Hyperbolic depth is thus defined as $\Lambda(X) = -\ln(1 - \sigma(X))$. This reveals that perfect geometric routing is an asymptotic boundary, bounded by universal regime transitions scaling multiplicatively by $e^2 \approx 7.389$. Notably, this identical scaling law and hyperbolic ladder describe the closure geometry of additive prime $k$-tuples, suggesting the divisor landscape and network routing optimization share a unified geometric constraint field.

\section{Empirical Validation: The Space Invariant Signal}
To isolate the mechanism generating geometric separation within this architecture, a pure measurement of the complex transport channel was executed using a train-free geometric probe. We utilized a cosine metric between matched and mismatched prototype class pairs across three candidate operative spaces:
\begin{enumerate}
    \item \textbf{CURRENT\_BASELINE:} $R^{16}$ hybrid space (12 angular + 4 magnitude dimensions).
    \item \textbf{R4\_FLAT:} $R^4$ space (compressed to the core harmonic block).
    \item \textbf{REDUCED\_BLOCK:} $R^{12}$ space (angular dimensions only, magnitude stripped).
\end{enumerate}

\begin{table}[h]
\centering
\begin{tabular}{lccc}
\toprule
\textbf{Space Type} & \textbf{Mean Matched Cos} & \textbf{Mean Mismatched Cos} & \textbf{Separation Delta} \\
\midrule
CURRENT\_BASELINE ($R^{16}$) & 1.000000 & -0.333333 & 1.333333 \\
R4\_FLAT ($R^4$)            & 1.000000 & -0.333333 & 1.333333 \\
REDUCED\_BLOCK ($R^{12}$)   & 1.000000 & -0.333333 & 1.333333 \\
\bottomrule
\end{tabular}
\caption{Empirical measurement of geometric separation signal across ambient space representations.}
\label{tab:signal_probe}
\end{table}

\subsection{Interpretation of Results}
The separation delta ($\Delta = 1.333333$) remains completely rigid across all spaces. The geometric signal is fully localized within the 2i-rotated harmonic pair encoding the classes as unit-norm vectors. Introducing or stripping ambient Euclidean dimensions (magnitude variance or unrelated angular blocks) does not alter the separation metric. This confirms that the observed routing advantages are not artifacts of high-dimensional ambient capacity, but are governed strictly by space-invariant, local phase-aligned structures. 

\section{Conclusion and Future Work}
The Riemannian-Complex architecture formalizes the intuitions derived from earlier Hopf-base discretizations. By explicitly modeling computation as phase-aware complex transport projected over a curved manifold, we enforce the structural invariance required for highly efficient routing. 

Future evaluations will focus on the stability of this architecture in dynamic, temporally extended environments. We propose that systems grounded in space-invariant geometric closure will natively maintain coherence across layered, consequential interaction loops—such as those found in multi-agent persistent state spaces—far more robustly than comparable Euclidean linear systems.

\bibliography{references}

\end{document}